% Vyžaduje balíček: graphicx. Obrázky: analyza/analyticka_cast/obrazky/
% (\graphicspath{{analyza/analyticka_cast/obrazky/}}).
\section{Typ a způsob využití objektů}

Další analýza se týká pouze aktuálního stavu, tedy 102\,220 objektů. Podle
typu je 60,8\,\% objektů budov s číslem popisným, 19,0\,\% budov s číslem
evidenčním a 20,2\,\% budov bez čísla popisného či evidenčního; typ objektu
je vyplněn u všech objektů.

Způsob využití je výrazně koncentrovaný (obrázek~\ref{fig:vyuziti}).
Téměř polovinu objektů tvoří rodinné domy (45,8\,\%), následují garáže
(17,4\,\%), stavby pro rodinnou rekreaci (11,0\,\%), bytové domy (7,6\,\%)
a jiné stavby (6,0\,\%). Každý z ostatních způsobů využití má podíl nižší
než 3\,\%. Pro přehlednost jsou sloučeny hodnoty „objekt občanské
vybavenosti“ a „stavba občanského vybavení“, z nichž první byla
v číselníku RÚIAN v roce 2016 nahrazena druhou, ale ve starších záznamech
zůstala.

\begin{figure}[htbp]
\centering
\includegraphics[width=\textwidth]{06_zpusob_vyuziti}
\caption{Způsob využití aktuálních stavebních objektů}
\label{fig:vyuziti}
\end{figure}

Struktura zástavby se liší mezi městem Brnem a okolními obcemi
(obrázek~\ref{fig:vyuziti-brno}). V okolních obcích dominují rodinné domy
(56,8\,\% oproti 41,9\,\% v Brně), zatímco v Brně je výrazně více
bytových domů (9,9\,\% oproti 1,7\,\%) a garáží (19,0\,\% oproti 13,8\,\%).
Podíl staveb pro rodinnou rekreaci se mezi oběma oblastmi téměř neliší.

\begin{figure}[htbp]
\centering
\includegraphics[width=\textwidth]{07_vyuziti_brno_okoli}
\caption{Nejčastější způsoby využití v Brně a v okolních obcích}
\label{fig:vyuziti-brno}
\end{figure}

\section{Podlažnost, velikost a počet bytů}

Počet podlaží je vyplněn u 66\,\% aktuálních objektů, zastavěná plocha
u 82\,\% a počet bytů u 83\,\%; všechny uvedené podíly a statistiky jsou
počítány z vyplněných hodnot (tabulka~\ref{tab:numericke}).

\begin{table}[htbp]
\centering
\caption{Souhrnné charakteristiky číselných atributů aktuálních objektů}
\label{tab:numericke}
\begin{tabular}{l|r|r|r}
Veličina & Počet hodnot & Medián & Maximum \tabularnewline \hline
Počet podlaží & 67\,805 & 2 & 32 \tabularnewline \hline
Zastavěná plocha [m$^2$] & 83\,296 & 105 & 95\,220 \tabularnewline \hline
Počet bytů (všechny objekty) & 84\,958 & 1 & 346 \tabularnewline \hline
Počet bytů (bytové domy) & 7\,775 & 15 & 346 \tabularnewline
\end{tabular}
\end{table}

\subsection{Podlažnost}

Zástavba je převážně nízkopodlažní (obrázek~\ref{fig:podlaznost}):
34\,\% objektů má jedno podlaží, 40\,\% dvě a 13\,\% tři. Šest a více podlaží
má jen asi 5\,\% objektů, nejvyšší objekt má 32 podlaží.

\begin{figure}[htbp]
\centering
\includegraphics[width=0.9\textwidth]{08_podlaznost}
\caption{Rozdělení objektů podle počtu podlaží}
\label{fig:podlaznost}
\end{figure}

Podlažnost úzce souvisí se způsobem využití
(obrázek~\ref{fig:podlaznost-vyuziti}). Rodinné domy mají nejčastěji dvě
podlaží (52\,\%), garáže jsou ve 94\,\% jednopodlažní a bytové domy mají
v 85\,\% čtyři a více podlaží (48\,\% čtyři až pět, 37\,\% šest a více).
Využití objektu je tak pro odhad jeho výšky silným ukazatelem, což je
užitečné zejména pro objekty s chybějící podlažností.

\begin{figure}[htbp]
\centering
\includegraphics[width=\textwidth]{09_podlaznost_podle_vyuziti}
\caption{Podlažnost podle způsobu využití (podíly objektů s vyplněnou podlažností)}
\label{fig:podlaznost-vyuziti}
\end{figure}

\subsection{Zastavěná plocha}

Medián zastavěné plochy je 105\,m$^2$ a polovina objektů leží mezi 59
a 166\,m$^2$; rozdělení má dlouhý pravý chvost, nejrozsáhlejší objekt zabírá
95\,220\,m$^2$. Histogram v logaritmickém měřítku
(obrázek~\ref{fig:plocha}) je dvouvrcholový: hlavní vrchol kolem 100\,m$^2$
odpovídá rodinným domům, menší vrchol kolem 20\,m$^2$ garážím a rekreačním
stavbám.

\begin{figure}[htbp]
\centering
\includegraphics[width=0.9\textwidth]{10_zastavena_plocha}
\caption{Rozdělení zastavěné plochy objektů (logaritmické měřítko)}
\label{fig:plocha}
\end{figure}

Rozdíly mezi způsoby využití jsou řádové (obrázek~\ref{fig:plocha-vyuziti}).
Medián činí 20\,m$^2$ u garáží, 29\,m$^2$ u rekreačních staveb, 110\,m$^2$
u rodinných domů, 298\,m$^2$ u bytových domů, 431\,m$^2$ u staveb
občanské vybavenosti a 692\,m$^2$ u staveb pro výrobu a skladování.

\begin{figure}[htbp]
\centering
\includegraphics[width=\textwidth]{11_plocha_podle_vyuziti}
\caption{Zastavěná plocha podle způsobu využití. Obdélník zahrnuje střední polovinu objektů (25.--75.\,percentil), oranžová čára je medián a vousy ukazují rozsah 5.--95.\,percentilu; n je počet objektů s vyplněnou plochou.}
\label{fig:plocha-vyuziti}
\end{figure}

\subsection{Počet bytů}

Počet bytů má smysl sledovat především u bytových domů
(obrázek~\ref{fig:byty}). Medián je 15 bytů, horní kvartil 25 bytů;
více než 60 bytů má asi 4\,\% domů a největší dům jich eviduje 346.

\begin{figure}[htbp]
\centering
\includegraphics[width=0.9\textwidth]{12_pocet_bytu_bytove_domy}
\caption{Počet bytů v bytových domech (hodnoty nad 60 jsou sloučeny)}
\label{fig:byty}
\end{figure}
